% ============================================================================
% report/experiments.tex — phần "Thiết lập thực nghiệm" (fragment, không có
% preamble). Dùng qua \input{} sau report/methodology.tex.
% Preamble cần: amsmath, booktabs, tabularx, array, xcolor.
% Các kết quả quan sát được nằm trong experimental_findings.tex; các caveat về
% provenance nằm trong reproducibility_limitations.tex.
% ============================================================================

\section{Thiết lập thực nghiệm}

Chương này trình bày điều kiện triển khai và phạm vi đánh giá của ba hệ thống
E1, E2 và E3-custom. Kiến trúc mô hình, quy trình huấn luyện, tiêu chí lựa chọn
checkpoint, chuẩn hoá đầu ra và phương pháp chấm điểm đã được trình bày trong
phần Phương pháp. Vì vậy, phần này chỉ mô tả phần cứng, môi trường phần mềm,
thiết lập ngẫu nhiên, các ràng buộc tài nguyên và cấu hình triển khai cần thiết
để diễn giải kết quả.

\subsection{Môi trường tính toán}

\subsubsection{Phần cứng}

Ba hệ thống được huấn luyện trên cùng một máy \texttt{laboratory}, sử dụng một
GPU \texttt{NVIDIA GeForce RTX 3060} với 12.488.343.552 byte bộ nhớ (xấp xỉ
12,5~GB). Các run manifest đều ghi nhận
\texttt{distributed\_world\_size=1}; do đó, không kết quả nào trong báo cáo
được tạo bởi thiết lập đa GPU.

E1 và E3-custom được chạy trên
\texttt{Linux-7.0.0-28-generic-x86\_64-with-glibc2.39}. Đoạn chạy tiếp hiện được
ghi trong manifest của E2 sử dụng
\texttt{Linux-7.0.0-30-generic-x86\_64-with-glibc2.39}. Sự khác biệt này phản
ánh trạng thái hệ thống tại thời điểm chạy, không phải khác biệt về phần cứng.
Các notebook Kaggle trong repo chỉ cung cấp đường chạy dự phòng và không sinh
ra các kết quả được báo cáo.

\subsubsection{Môi trường phần mềm}

Dự án duy trì ba môi trường Conda tách biệt để cô lập các phụ thuộc của
Fairseq, QLoRA và bộ đánh giá. Cách tổ chức này tránh việc ép Fairseq 0.12.2 và
các kernel lượng tử hoá 4-bit của \texttt{bitsandbytes} vào cùng một software
stack. Bảng~\ref{tab:compute-env} ghi phiên bản thực tế trong các artifact hiện
có; hàng E2 tương ứng với đoạn chạy tiếp đến step 910.

\begin{table}[tbp]
  \centering
  \small
  \renewcommand{\arraystretch}{1.18}
  \caption{Các môi trường Conda và phiên bản được ghi trong artifact hiện tại.}
  \label{tab:compute-env}
  \begin{tabularx}{\textwidth}{@{}>{\ttfamily}p{2.3cm}p{2.3cm}p{3.15cm}X@{}}
    \toprule
    \normalfont\textbf{Môi trường} & \textbf{Dùng cho}
      & \textbf{Python / PyTorch / CUDA} & \textbf{Thư viện chính} \\
    \midrule
    .conda/\allowbreak fairseq
      & E1, E3-custom
      & 3.10.20 / 1.13.1 / 11.7
      & \texttt{fairseq} 0.12.2 \\
    .conda/llm
      & E2
      & 3.10.20 / 2.6.0+cu124 / 12.4
      & \texttt{transformers} 4.57.6; \texttt{peft} 0.20.0;
        \texttt{bitsandbytes} 0.50.1; \texttt{accelerate} 1.14.0;
        \texttt{datasets} 4.8.5 \\
    .conda/eval
      & Chấm điểm, kiểm toán, kiểm thử
      & 3.10.20 / --- / ---
      & \texttt{sacrebleu} 2.6.0; \texttt{PyYAML} 6.0.2;
        \texttt{pytest} 8.4.1 \\
    \bottomrule
  \end{tabularx}
\end{table}

Tệp khai báo môi trường E2 từng đặt PyTorch 2.5.1, trong khi manifest của đoạn
chạy tiếp ghi \texttt{2.6.0+cu124}; Bảng~\ref{tab:compute-env} ưu tiên phiên bản
đã được ghi nhận khi chạy. Việc chấm điểm được tách khỏi các môi trường huấn
luyện: mọi metric artifact dùng SacreBLEU 2.6.0 và lưu chữ ký chứa
\texttt{version:2.6.0}. Vì vậy, cả ba hệ thống được chấm bằng cùng một phiên bản
bộ đánh giá, dù software stack dùng để huấn luyện khác nhau.

\subsubsection{Kiểm soát ngẫu nhiên và khả năng lặp lại}

Cả ba cấu hình đặt \texttt{seed: 42}. Với Fairseq, pipeline:
\begin{itemize}
  \item bật \texttt{cudnn.deterministic} và tắt
  \texttt{cudnn.benchmark};
  \item gọi \texttt{use\_deterministic\_algorithms} ở chế độ
  \texttt{warn\_only};
  \item đặt \texttt{CUBLAS\_WORKSPACE\_CONFIG=:4096:8} và
  \texttt{PYTHONHASHSEED} cho tiến trình huấn luyện.
\end{itemize}
Với QLoRA, \texttt{set\_seed} được gọi trước khi huấn luyện và sinh dự đoán;
\texttt{TrainingArguments} nhận cả \texttt{seed} và \texttt{data\_seed}; bước
packing sử dụng một generator riêng cũng được khởi tạo từ 42.

Các thiết lập trên làm giảm biến thiên và tăng khả năng lặp lại khi giữ nguyên
ngăn xếp phần cứng và phần mềm. Tuy nhiên, chế độ \texttt{warn\_only} không buộc mọi
CUDA operation phải tất định, nên chúng không cấu thành cam kết tái lập
bit-for-bit khi đổi GPU, driver, CUDA hoặc phiên bản thư viện.

\subsection{Cấu hình triển khai}

\subsubsection{Ràng buộc tài nguyên}

Giới hạn 12,5~GB VRAM chi phối một số lựa chọn triển khai. Với E2, trọng số nền
Qwen3-8B được lượng tử hoá NF4 kèm double quantization; chỉ LoRA adapter nhận
gradient. Cấu hình đồng thời sử dụng gradient checkpointing,
\texttt{paged\_adamw\_8bit}, một packed window trên mỗi thiết bị, tích luỹ
gradient qua 16 bước và độ dài cửa sổ 768 token. Với E1 và E3-custom,
\texttt{max\_tokens\_per\_gpu=1000} kết hợp \texttt{update\_freq=4}, tương ứng
với ngân sách khai báo khoảng 4.000 token cho mỗi lần cập nhật tham số.

Các lựa chọn này mô tả cách ba hệ thống được triển khai trên phần cứng sẵn có;
chúng không phải kết quả của một nghiên cứu ablation về bộ nhớ hoặc batch size.
Ảnh hưởng quan sát được của giới hạn này được trình bày cùng các phát hiện thực
nghiệm.

\subsubsection{Tóm tắt cấu hình thực tế}

Bảng~\ref{tab:exp-recap} cung cấp một bản tóm tắt để đối chiếu nhanh; định nghĩa
đầy đủ của từng mô hình và realized protocol nằm trong phần Phương pháp.

\begin{table}[tbp]
  \centering
  \footnotesize
  \renewcommand{\arraystretch}{1.16}
  \setlength{\tabcolsep}{3.5pt}
  \caption{Tóm tắt cấu hình triển khai của ba hệ thống được đánh giá.}
  \label{tab:exp-recap}
  \begin{tabularx}{\textwidth}{@{}>{\bfseries}p{2.35cm}XXX@{}}
    \toprule
    & E1 & E2 & E3-custom \\
    \midrule
    Khung mô hình
      & Transformer 6+6, $d=512$, 67,65M tham số
      & Qwen3-8B đóng băng + LoRA $r=16$, $\alpha=32$
      & Như E1, nguồn được bổ sung gợi ý, 74,64M tham số \\
    Độ chính xác
      & fp16
      & Trọng số nền NF4 4-bit, double quantization; tính toán fp16
      & fp16 \\
    Bộ tối ưu
      & Adam, $5\times10^{-4}$, inverse-sqrt, warmup 8.000
      & Paged AdamW 8-bit, $2\times10^{-4}$, cosine, warmup ratio 0,05
      & Như E1 \\
    Đơn vị batch
      & 1.000 token $\times$ tích luỹ 4
      & 1 packed window $\times$ tích luỹ 16
      & Như E1 \\
    Ngân sách thực tế
      & Dừng ở epoch 171 / update 24.791
      & Chạy tiếp từ 3 epoch / 546 step đến 5 epoch / 910 step
      & Chạm 50.000 update ở epoch 87 \\
    Chọn checkpoint
      & Validation BLEU, epoch 151
      & Validation \texttt{eval\_loss}, step 400
      & Validation BLEU, epoch 77 \\
    Giải mã
      & Beam 7; $1{,}2|x|+10$
      & Beam 4; $\min(512,\lceil2|x|\rceil+32)$
      & Như E1 \\
    \bottomrule
  \end{tabularx}
\end{table}

Đơn vị batch của hai backend không quy đổi trực tiếp: Fairseq giới hạn số token
trên mỗi GPU, còn E2 đếm các cửa sổ đã packing. Do đó, hàng “Đơn vị batch”
không được diễn giải như một so sánh về số câu trong mỗi update.

\subsection{Phạm vi thực nghiệm}

Báo cáo so sánh ba hệ thống dịch Việt $\rightarrow$ Hán văn trên cùng các phân
hoạch cố định của ngữ liệu \emph{Đại Việt sử ký toàn thư}: E1 là Transformer
huấn luyện từ đầu; E2 là Qwen3-8B thích nghi bằng QLoRA; E3-custom là E1 với
câu nguồn được bổ sung gợi ý Hán tự từ kho tri thức. Đây là so sánh cấp hệ
thống, không phải một chuỗi ablation kiểm soát từng biến. Các bất biến về dữ
liệu, quy trình lựa chọn trên validation, cổng đóng băng tập test và bộ đánh
giá chung đã được định nghĩa trong phần Phương pháp.

Hai hệ thống trong đặc tả ban đầu không thuộc tập kết quả. E3 chính thức phụ
thuộc integration \texttt{knowledge\_v2} chưa có trong repo và vì vậy ở trạng
thái blocked; hệ thống đã chạy phải luôn được ghi là E3-custom. E4, tức Qwen3
kết hợp cơ chế truy hồi của E3-custom, mới được cài đặt và chạy preflight ước
lượng chi phí, chưa được huấn luyện. Báo cáo không gán kết quả định lượng cho
E3 chính thức hoặc E4.
